% Ad Familiares 11.29 (ad-familiares-11-29)
% Source: https://raw.githubusercontent.com/PerseusDL/canonical-latinLit/master/data/phi0474/phi056/phi0474.phi056.perseus-lat1.xml
% Fetched by scripts/fetch_latin.py

% Dateline (Perseus): Scr. vel ex. m. Iun. vel in. Quint. a. 710(44). .

\ciceroLetterOpener{CICERO OPPIO S. D.}

\ciceroSection{1} dubitanti mihi, quod scit Atticus noster, de hoc toto consilio profectionis, quod in utramque partem in mentem multa veniebant, magnum pondus accessit ad tollendam dubitationem iudicium et consilium tuum ; nam et scripsisti aperte quid tibi videretur et Atticus ad me sermonem tuum pertulit. semper iudicavi in te et in capiendo consilio prudentiam summam esse et in dando fidem maximeque sum expertus, cum initio civilis belli per litteras te consuluissem quid mihi faciendum esse censeres, eundumne ad Pompeium an manendum in Italia. suasisti ut consulerem dignitati meae ; ex quo quid sentires intellexi et sum admiratus fidem tuam et in consilio dando religionem, quod, cum aliud malle amicissimum tuum putares, antiquius tibi officium meum quam illius voluntas fuit.

\ciceroSection{2} equidem et ante hoc tempus te dilexi et semper me a te diligi sensi et, cum abessem atque in magnis periculis essem, et me absentem et meos praesentis a te cultos et defensos esse memini et, post meum reditum quam familiariter mecum vixeris quaeque ego de te et senserim et praedicarim, omnis qui solent haec animadvertere testis habemus. gravissimum vero iudicium de mea fide et constantia fecisti, cum post mortem Caesaris totum te ad amicitiam meam contulisti ; quod tuum iudicium nisi mea summa bene volentia erga te omnibusque meritis comprobaro, ipse me hominem non putabo.

\ciceroSection{3} tu , mi Oppi, conservabis amorem tuum (etsi more magis hoc quidem scribo quam quo te admonendum putem) meaque omnia tuebere ; quae tibi ne ignota essent Attico mandavi ; a me autem, cum paulum oti nacti erimus, uberiores litteras exspectato. da operam ut valeas. hoc mihi gratius facere nihil potes.
